\section*{Hoofdstuk \ref{ch:g9:atoms-to-galaxies} --- Van atomen tot sterrenstelsels: schalen van het heelal}

\begin{solution}{exo:g9:atoms-to-galaxies:1}
\omterm{def:g9:atoms-to-galaxies:atom}{Kern} ($10^{-15}$ \unit{m}), \omterm{def:g9:atoms-to-galaxies:atom}{atoom} ($10^{-10}$),
\omterm{def:g7:states-of-matter:molecule}{molecuul} (rond $10^{-9}$), mier ($10^{-2}$, of $10^{-3}$ voor
een bescheiden exemplaar).
\end{solution}

\begin{solution}{exo:g9:atoms-to-galaxies:2}
Een minuscule, zware, positief geladen \omterm{def:g9:atoms-to-galaxies:atom}{kern} in het midden; een
wolk van vrijwel gewichtloze negatieve \omterm{def:g9:atoms-to-galaxies:atom}{elektronen} ver
daarbuiten. De afmetingen: \omterm{def:g9:atoms-to-galaxies:atom}{atoom} rond $10^{-10}$ \unit{m},
\omterm{def:g9:atoms-to-galaxies:atom}{kern} rond $10^{-15}$ --- de kloof van honderdduizend die de
erwt-in-een-stadion oplevert.
\end{solution}

\begin{solution}{exo:g9:atoms-to-galaxies:3}
\omterm{def:g9:atoms-to-galaxies:atom}{Elektronen} --- negatief, marcherend van $-$ naar $+$:
tegengesteld aan de afspraak. De wetten overleven omdat ze nooit
afhingen van wie er marcheert: stromen,
knooppuntsommen, de portretten van Ohm en de
vermogensrekeningen lezen precies hetzelfde met de pijl die
wereldwijd is afgesproken.
\end{solution}

\begin{solution}{exo:g9:atoms-to-galaxies:4}
De dichtstbijzijnde macht van tien. Lengte: $10^{0}$ \unit{m};
klaslokaal: $10^{1}$; diameter van de Aarde: $10^{7}$.
\end{solution}

\begin{solution}{exo:g9:atoms-to-galaxies:5}
\omterm{def:g9:atoms-to-galaxies:atom}{Atoom} tot \omterm{def:g9:atoms-to-galaxies:atom}{kern}: vijf sporten. Jij tot de Aarde:
zeven. Afstand tot de Zon ($10^{11}$) tot de dichtstbijzijnde
\omterm{def:g4:solar-system:star}{ster} ($4 \times 10^{16}$): vijf tot zes sporten --- een paar
honderdduizend keer.
\end{solution}

\begin{solution}{exo:g9:atoms-to-galaxies:6}
Het \omterm{def:g9:atoms-to-galaxies:atom}{atoom} is $10^{-10}$ \unit{m} vrijwel niets rond een
\omterm{def:g9:atoms-to-galaxies:atom}{kern} van $10^{-15}$: pak \omterm{def:g9:atoms-to-galaxies:atom}{atomen} samen tot een
tafel en de leegte gaat mee. Vastheid is de elektrische weigering van de
elektronenwolken om te overlappen --- leegte, stevig verdedigd.
\end{solution}

\begin{solution}{exo:g9:atoms-to-galaxies:7}
Millimeter, $10^{-3}$; \omterm{def:g9:atoms-to-galaxies:atom}{atoom}, $10^{-10}$: zeven sporten ---
ongeveer tien miljoen \omterm{def:g9:atoms-to-galaxies:atom}{atomen} in de rij.
\end{solution}

\begin{solution}{exo:g9:atoms-to-galaxies:8}
\omterm{def:g8:speed-of-light:lightyear}{Lichtjaar}: $10^{16}$ \unit{m} (negenenhalf duizend miljoen
miljoen). \omterm{ex:g5:stars-night-sky:milkyway}{Melkweg}: $10^{21}$. Diepst in kaart gebrachte hemel:
rond $10^{26}$.
\end{solution}

\begin{solution}{exo:g9:atoms-to-galaxies:9}
$10^{5}$ erwten: $0.01 \times 10^{5} = \qty{1000}{m}$ --- een
``stadion'' van een kilometer breed: groter dan welk echt stadion ook,
maar de vergelijking houdt stand --- een erwt op de middenstip, en de
volgende erwten een kilometer verderop.
\end{solution}

\begin{solution}{exo:g9:atoms-to-galaxies:10}
Druppels in de zeeën: ongeveer $10^{21} \times 10^{5} = 10^{26}$ ---
dus de oude bewering reikte, letterlijk genomen, te ver: de zeeën
bevatten meer druppels dan een druppel \omterm{def:g7:states-of-matter:molecule}{moleculen} bevat, met
enkele sporten verschil. De eerlijke versie --- dat de
\omterm{def:g7:states-of-matter:molecule}{moleculen} in een druppel talrijker zijn dan wat dan ook dat je
in een mensenleven kunt tellen --- overleeft; ladders houden zelfs
geliefde beweringen eerlijk.
\end{solution}

\begin{solution}{exo:g9:atoms-to-galaxies:11}
Ongeveer tien sporten omlaag tot het \omterm{def:g9:atoms-to-galaxies:atom}{atoom} ($10^{0}$ tot
$10^{-10}$) en zesentwintig omhoog tot de diepe hemel --- dicht bij het
midden, een beetje naar het kleine gekanteld. Negen jaar onderwijs
brengen een lezer op eerlijke afstand van beide uiteinden van alles wat
gemeten is.
\end{solution}

\begin{solution}{exo:g9:atoms-to-galaxies:12}
De antwoorden zijn persoonlijk --- de oefening is de brief zelf. (Eén
voorbeeld, voor de vorm: aan de lezer uit groep drie van het hoofdstuk
over \omterm{def:g1:light-and-shadows:shadow}{schaduwen} --- ``de donkere tweelingbroer op de stoep zal
op een dag \omterm{def:g7:shadows-eclipses:solar}{zonsverduisteringen} verklaren''; aan de lezer uit groep zeven
van het energiehoofdstuk --- ``de ketens die je natrekt worden
formules met \omterm{def:g9:kinetic-energy-safety:joule}{joule} erin''; aan de lezer uit klas twee van de
wet van Ohm --- ``de portretten die je tekent zijn de gewoonte om wetten
te toetsen, en dat is het hele geheim''.)
\end{solution}

\begin{solution}{pb:g9:atoms-to-galaxies:1}
\textbf{1.} Schoolplein: \omterm{def:g2:measuring-time:second}{seconde} $2$; de Aarde: \omterm{def:g2:measuring-time:second}{seconde} $7$; de
\omterm{def:g6:describing-motion:trajectory}{baan} van de \omterm{def:g2:moon-in-the-sky:moon}{Maan}: \omterm{def:g2:measuring-time:second}{seconde} $9$.
\textbf{2.} \omterm{def:g2:measuring-time:second}{Seconde} $11$ (bij $1.5 \times 10^{11}$ \unit{m}); de
verteller: ``de zonneschijn op het kleed is acht minuten oud.''
\textbf{3.} Ruwweg \omterm{def:g2:measuring-time:second}{seconde} $13$ tot $16$: het kader bevat het hele
\omterm{def:g4:solar-system:family}{zonnestelsel} als een stip en nog geen enkele
\omterm{def:g4:solar-system:star}{ster} --- eerlijk gezegd bijna zwarte leegte, \omterm{def:g2:measuring-time:second}{seconden} lang: de
waarachtigste scènes van de film.
\textbf{4.} De diepst in kaart gebrachte hemel, $10^{26}$ \unit{m} aan
sterrenstelsels --- met het oudste \omterm{def:g1:light-and-shadows:source}{licht} erin ouder dan de Aarde
zelf.
\textbf{5.} Mier: \omterm{def:g2:measuring-time:second}{seconde} $2$ naar binnen; cel: \omterm{def:g2:measuring-time:second}{seconde} $5$;
\omterm{def:g7:states-of-matter:molecule}{molecuul}: \omterm{def:g2:measuring-time:second}{seconde} $9$; \omterm{def:g9:atoms-to-galaxies:atom}{atoom}: \omterm{def:g2:measuring-time:second}{seconde} $10$.
\textbf{6.} Door de eigen leegte van het \omterm{def:g9:atoms-to-galaxies:atom}{atoom} --- vijf \omterm{def:g2:measuring-time:second}{seconden}
niets tussen wolk en middelpunt: ``als dit \omterm{def:g9:atoms-to-galaxies:atom}{atoom} een stadion was,
hebben we de hoogste tribunes verlaten en vliegen we nog steeds naar een
erwt toe.''
\textbf{7.} De \omterm{def:g9:atoms-to-galaxies:atom}{kern} --- en eerlijk: ``hier eindigt onze ladder,
niet omdat de natuur ophoudt, maar omdat dit de kleinste sporten zijn
die mensen hebben gemeten; er wordt verder gebouwd.''
\textbf{8.} Naar buiten $26$ \omterm{def:g2:measuring-time:second}{seconden}; naar binnen $15$: een heelal van
eenenveertig \omterm{def:g2:measuring-time:second}{seconden}.
\textbf{9.} Elke \omterm{def:g2:measuring-time:second}{seconde} vermenigvuldigt het kader met tien --- gelijke
\emph{verhoudingen}, geen gelijke optellingen: de logica van de ladder
is dat een sport een factor is, en in factoren zijn de straatstap en de
sterrenhoopstap even groot.
\textbf{10.} Bijvoorbeeld: ``Hoe vol het er ook uitziet, bijna alles
hier is leeg --- materie en hemel zijn allebei zeldzame eilanden in een
wijde leegte.''
\textbf{11.} De microscoop (het lenzenhoofdstuk: twee
\omterm{def:g8:lenses-images:families}{convergerende lenzen} boven het kleine) en de telescoop (hetzelfde
hoofdstuk, dat het zwakke verre verzamelt) --- de sleutels van de twee
richtingen, allebei geslepen uit de optica van dit boek.
\textbf{12.} Bijvoorbeeld: ``Negen jaar, één ladder: van de telbare
pootjes van een lieveheersbeestje tot sterrenstelsels ouder dan de grond
onder je voeten, en elke sport werd gezet door metingen die iedereen mag
nagaan. Volgend jaar beginnen de vragen die we open lieten ---
\omterm{def:g2:pushes-and-pulls:force}{kracht} en beweging, de aard van het \omterm{def:g1:light-and-shadows:source}{licht}, het
hart van het \omterm{def:g9:atoms-to-galaxies:atom}{atoom} --- zich te beantwoorden; neem de ladder
mee.''
\end{solution}
